\scanpage{245}
\mspage{423}

\margin{one should \unsure{give its} \uline{formulas} in a precise \unsure{way}, \unsure{take} \illegible{2}!!!}

Now according to the general principles (not yet established, to tell the truth, in the profinite setting) the commutant of $\rho$ in the group in question is reduced to the group generated by $\rho$, $\{1,\rho,\rho^2\}$, so that the relation \unsure{(*)} takes the form
\[
\tau_\infty^{\prime\alpha_3}\widehat{\varphi}(f_\infty)\alpha=\rho^j
\qquad\text{for }j\in\mathbb Z/3\mathbb Z\ \unsure{\text{suitable}}
\quad\text{(well determined)}.
\]
We make it explicit by \uline{using} the value of $f_\infty$
\[
f_\infty^{-1}=g\rho(g)\rho^2(f_0^{-1})
=g\rho(g)=\rho^2(g)^{-1}=\rho^{-1}(\xi^{-1}),
\]
\inserted{N.B. here $\left\{\begin{gathered}f_0=f=1\\g=\xi=\alpha\rho(\alpha)^{-1}\end{gathered}\right.$}
thus the relation \deleted{\illegible{2}} becomes
\[
\tau_\infty^{\prime\alpha_3}\widehat{\varphi}\bigl(\rho^{-1}(\xi)\bigr)\alpha=\rho^j
\]
i.e., \deleted{by applying $\tau_\infty^{\prime\alpha_3}$ \unsure{to the}}
\[
\widehat{\varphi}\bigl(\rho^{-1}(\xi)\bigr)
=\underbrace{\tau_\infty^{\prime\alpha_3}\rho^j}_{\underbrace{\tau_\infty^{\prime\alpha_3}(\rho^j)}\tau_\infty^{\prime\alpha_3}}\alpha^{-1}
\LMcancelledtag{\deleted{\unsure{44|43}}}
\]
or also
\deleted{\[
\widehat{\varphi}\bigl(\rho^{-1}(\xi^{-1})\bigr)=\alpha\rho^{-j}\tau_\infty^{\prime\alpha_3}.
\tag{\unsure{44 bis|43 bis}}
\]}
since $\tau'_\infty(\rho)=\rho^{-1}$ thus
$\tau_\infty^{\prime\alpha_3}(\rho^j)=\rho^{(-1)^{\alpha_3}j}=\rho^{j'}$
\[
\widehat{\varphi}\bigl(\rho^{-1}(\xi)\bigr)=\rho^{j'}\tau_\infty^{\prime\alpha_3}\alpha^{-1}
\tag{44}
\]
where $\alpha_3=\alpha_3(\mu)\in\mathbb Z/2\mathbb Z$ defined by $(-1)^{\alpha_3}\equiv\mu\ (3)$ and $j'\in\mathbb Z/3\mathbb Z$ (well determined).

We shall transform this relation a little further, by writing
\[
\tau_\infty^{\prime\alpha_3}=(\sigma_\infty\tau_\infty)^{\alpha_3}=\sigma_\infty^{\alpha_3}\tau_\infty^{\alpha_3}
\qquad\text{whence}
\]
\[
\tau_\infty^{\prime\alpha_3}\alpha^{-1}=\sigma_\infty^{\alpha_3}\bigl(\underbrace{\alpha\tau_\infty^{\alpha_3}}_{\widetilde\alpha}\bigr)^{-1}
\tag{\unsure{*}}
\]
\mspage{244}

\scanpage{246}

where \deleted{$\unsure{\alpha\tau_\infty|\alpha_0}$}
\[
\widetilde\alpha=\alpha\tau_\infty^{\alpha_3}\deleted{\unsure{(\alpha_3)}}
\tag{45}
\]
has the advantage over $\alpha$ itself of being in $\widehat\Pi_{0,3}$ in all cases (\uline{N.B.} we have
$\alpha\in\widehat\Pi_{0,3}\cdot\{1,\tau_\infty\}$, and $\alpha\in\widehat\Pi_{0,3}$ iff $\alpha_3=0$).
It is equal to $\alpha$ iff $\alpha_3=0$ and in the case $\alpha_3=1$ \deleted{\illegible{1}} is directly related to $\xi$ by the formula
\[
\xi=\widetilde\alpha\ell_1^{-1}\rho(\widetilde\alpha)^{-1}
\tag{46}
\]
(deduced from (45), from $\xi=\alpha\rho(\alpha)^{-1}$, and from
\[
\tau_\infty\rho(\tau_\infty)^{-1}=\ell_1^{-1}
\tag{**}
\]
).
Thus (44) takes the definitive form
\[
\boxed{\widehat{\varphi}\bigl(\rho^{-1}(\xi)\bigr)=\rho^{j'}\sigma_\infty^{\alpha_3}\unsure{\lbrace}\widetilde\alpha^{-1}}
\qquad(\widetilde\alpha=\alpha\tau_\infty^{\alpha_3})
\tag{47}
\]
\margin{all the factors of this formula are \unsure{defined} in $\unsure{\widehat{\mathcal G}_{0,3}^{+}|\widetilde{\mathcal G}_{0,3}^{+}}$, i.e. of multiplier $+1$}
which can \unsure{also|thus} \unsure{be written|describe}, on passing to inverses
\[
\widehat{\varphi}\bigl(\rho^{-1}(\xi^{-1})\bigr)=\widetilde\alpha\sigma_\infty^{\alpha_3}\rho^{-j'}
\tag{47 bis}
\]
or again, solving for $\widetilde\alpha$
\[
\underset{\substack{\Vert\\\alpha\tau_\infty^{\alpha_3}}}{\widetilde\alpha}
=\widehat{\varphi}\bigl(\rho^{-1}(\xi^{-1})\bigr)\rho_{\deleted{\unsure{\infty}}}^{j'}\sigma_\infty^{\alpha_3}.
\tag{48}
\]
We shall compare this formula with formula (10), which is solved for $\beta$ in the form
\[
\widehat{\varphi}(\xi)=\beta\sigma_\infty^i
\tag*{(49)=(10)}
\]
\[
\beta=\widehat{\varphi}(\xi)\sigma_\infty^i
\tag{50}
\]
\mspage{245}

\scanpage{247}
\mspage{424}

The fact that formula (47) is indeed equivalent to (10), or again to (49), when $\alpha\in\unsure{\widehat{\widetilde\Pi}_{0,3}|\widetilde{\widehat\Pi}_{0,3}}$ and $\beta\in\widehat\Pi_{0,3}$ are related by $\xi=\eta\ell_1^\nu$ (where $\xi=\alpha\rho(\alpha)^{-1}$, $\eta=\beta\sigma_\infty(\unsure{\eta})^{-1}$), would deserve a direct verification. The interest of (47), or of \deleted{(48)} its form solved for $\alpha$ (48), just as of (10)\inserted{$=(49)$} or (50), consists in the fact that these formulas give the ``delicate'' parameters\inserted{$\alpha,\beta$}, in terms of the ``coarse parameter'' $\xi$. Both are formulas in $\unsure{\widehat{\mathcal G}_{0,3}^{+}}$ --- on reducing modulo $\widehat\Pi_{0,3}$, they give formulas in $\unsure{\mathfrak S_3|\mathcal S_3}$, namely (taking into account that $\widetilde\alpha,\beta\in\widehat\Pi_{0,3}$):
\deleted{\unsure{i.e.}}
\[
\left\{
\begin{aligned}
\widehat{\varphi}^{\cdot}\bigl(\rho^{-1}(\xi)\bigr)&=\dot\rho^{j'}\dot\sigma_\infty^{\alpha_3}\\
\widehat{\varphi}^{\cdot}(\xi)&=\dot\sigma_\infty^i
\end{aligned}
\right.
\tag{51}
\]
which \unsure{gives back|gives} exactly \deleted{\unsure{the}}, in terms of the elements $\widehat{\varphi}^{\cdot}(\xi)$ and $\widehat{\varphi}^{\cdot}\bigl(\rho^{-1}(\xi)\bigr)$ of $\unsure{\mathfrak S_3|\mathcal S_3}$, the values of $i,\alpha_3\in\mathbb Z/2\mathbb Z$ and of $j'\in\mathbb Z/3\mathbb Z$. (We should see whether $i,j'$ can, like $\alpha_3$, be made explicit in terms of the multiplier $\mu$, or even in terms of its
\margin{and also see whether $(i,\alpha_3,j')$ can take \unsure{all sorts} of values \unsure{in} $\mathbb Z/2\times\mathbb Z/2\times\mathbb Z/3$}
\mspage{246}

\scanpage{248}

quadratic residues $\varepsilon_\ell(\mu)$) We can supplement them with a formula deduced from one and \deleted{from} the other, taking into account that
\[
\xi\rho(\xi)\underbrace{\rho^2(\xi)}_{\rho^{-1}(\xi)}=1
\qquad\text{i.e.}\qquad
\begin{gathered}
\deleted{\rho^{-1}(\xi^{-1})=\xi\rho(\xi)}\\
\rho^{-1}(\xi)\xi=\rho(\xi)^{-1}
\end{gathered}
\]
we therefore find, multiplying (47) and (10)
\[
\widehat{\varphi}\bigl(\underbrace{\rho^{-1}(\xi)\cdot\xi}_{\rho(\xi)^{-1}}\bigr)
=\rho^{j'}\sigma_\infty^{\alpha_3}\widetilde\alpha^{-1}\beta\sigma_\infty^i
\]
i.e. again
\[
\widehat{\varphi}\bigl(\rho(\xi)\bigr)
=\sigma_\infty^i(\beta^{-1}\widetilde\alpha)\sigma_\infty^{\alpha_3}\rho^{-j'}
\tag{52}
\]
i.e., in the form solved for $\beta^{-1}\widetilde\alpha$
\[
\underbrace{\beta^{-1}\widetilde\alpha}_{(\beta^{-1}\alpha)\tau_\infty^{\alpha_3}}
=\sigma_\infty^i\widehat{\varphi}\bigl(\rho(\xi)\bigr)\sigma_\infty^{\alpha_3}\rho^{j'}.
\tag{53}
\]
We can consider that the three main ``delicate parameters'' are the quantities $\beta$, $\alpha$ and $\beta^{-1}\alpha$, which are expressed directly, in a remarkably simple way, in terms of the values of $\widehat{\varphi}$ for $\xi$, $\rho^{-1}(\xi)$ and $\unsure{\rho^{-1}|\rho^{-2}}(\xi)=\rho(\xi)$, which would dispense \deleted{\illegible{1}} probably (by
\mspage{247}

\scanpage{249}
\mspage{425}

carrying out the calculations in terms of $\xi,\eta$ exclusively, which would require a little more effort...) with appealing to delicate results \unsure{not} yet proved, \inserted{of the type $\xi=\alpha\rho(\alpha)^{-1}$, $\eta=\beta\sigma_\infty(\beta)^{-1}$,} to justify the introduction of these parameters $\alpha,\beta$.
\margin{(recapitulation)}
\[
\left\{
\begin{aligned}
\widehat{\varphi}(\xi)&=\beta\sigma_\infty^i\\
\widehat{\varphi}\bigl(\rho^{-1}(\xi)\bigr)&=\rho^{j'}\sigma_\infty^{\alpha_3}\widetilde\alpha^{-1}
\qquad(\widetilde\alpha=\alpha\tau_\infty^{\alpha_3})\\
\widehat{\varphi}\bigl(\underbrace{\rho^{\unsure{-2|-1}}(\xi)}_{\rho(\xi)}\bigr)
&=\underbrace{\sigma_\infty^i(\beta^{-1}\widetilde\alpha)\sigma_\infty^{\alpha_3}}_{\sigma_\infty^i(\beta^{-1}\widetilde\alpha)\sigma_\infty^{\alpha_3+i}\rho^{-j'}}\rho^{-j'}
\end{aligned}
\right.
\tag{54}
\]

Suppose for simplicity
\[
\alpha_3=0\quad\text{i.e.}\quad\alpha=\widetilde\alpha\quad\text{i.e.}\quad\alpha\in\widehat\Pi_{0,3}
\tag{55}
\]
Then we obtain from the second formula (54), taking inverses
\[
\widehat{\varphi}\bigl(\rho^{-1}(\xi^{-1})\bigr)=\alpha\rho^{-j'}
\]
and the $x\rho(x)^{-1}$ of both sides
\[
\xi=\widehat{\varphi}\bigl(\rho^{-1}(\xi^{-1})\bigr)\rho\Bigl(\widehat{\varphi}\bigl(\rho^{-1}(\xi^{-1})\bigr)\Bigr)^{-1}
\tag{56}
\]
which already implies the first formula (14) in the fundamental formulas (14) (15).
\mspage{248}

\scanpage{250}
\mspage{425}

We thus find in this case ($\mu\equiv1\ (3)$ i.e.
$\nu\equiv0\ (3)$) that the two relations (14) (15)
are also equivalent to \unsure{the set|one} of the two relations
\margin{Valid only if $\nu\equiv0\ (3)$ i.e.
$\mu\equiv1\ (3)$ i.e. $\varepsilon_3(\mu)=1$.}
\[
\begin{array}{c@{\quad}c}
\begin{gathered}(57)\\[1.5ex](58)\end{gathered}&
\boxed{\begin{aligned}
\xi&=
\widehat\varphi\bigl(\rho^{-1}(\xi^{-1})\bigr)
\rho\bigl(\widehat\varphi(\rho^{-1}(\xi^{-1}))\bigr)^{-1}
\\
\xi&=\widehat\varphi(\xi)\,
\sigma_\infty\bigl(\widehat\varphi(\xi)\bigr)^{-1}
\ell_1^\nu
\end{aligned}}
\end{array}
\]
\inserted{with $\nu\in\widehat{\mathbb Z}$,
$\nu\equiv0\ (3)$}
which therefore characterize the
$\unsure{u_0|u_{\cdot}}=u_{\xi}\in
\unsure{\mathcal M_{0,3}^{!\sim}|\mathcal M_{0,3}^{!\approx}}(0)$,
whose multiplier $\mu$ is $\equiv1\ (3)$.

\uline{Remark} The only element of
$\mathcal M_{0,3}^{!\approx}(0)$
\inserted{$\simeq\mathcal N_{0,3}^{!\approx}$}
that I know how to make explicit, apart from the unit element
(and \unsure{one} even in
$\mathcal N_{0,3}^{!\sim}\supset
\unsure{\mathcal N_{4,3}^{!\approx}|\mathcal N_{g,3}^{!\approx}}$)
is \deleted{\illegible{1}}\inserted{the unique} element
$\tau\in\mathcal T_{0,3}^{!}$
\inserted{$\subset\mathcal N_{0,3}^{!\approx}$},
which in $\mathcal M_{0,3}^{!\approx}(0)$ corresponds
\inserted{to the element $\tau_\infty$, \deleted{\illegible{1}}
$=u_{\xi}=u_{\xi,\eta}$}
\deleted{to the invariants}
\[
\begin{aligned}
(59)\qquad & h=\beta_0=1,\quad f=1,\quad
\alpha=\alpha_0=\tau_\infty,\quad
\xi=g=\ell_1^{-1},\quad\eta=1
\\
&\nu=-1,\quad\mu=-1,
\end{aligned}
\]
\margin{Careful, here $\alpha_3=1$ so (57) is not
\unsure{valid|applicable}.}
\deleted{which}\inserted{\unsure{and} $\xi,\nu$}
do satisfy \deleted{\unsure{(57) (58)}}
\uline{$\xi\rho(\xi)
\rho^2(\xi)=1$}
\inserted{i.e. $\ell_0^{-1}\ell_1^{-1}\ell_\infty^{-1}=1$}
\margin{and the relations (54) \unsure{page above}}
\deleted{
\[
\widehat\varphi(\xi)=\ell_0,\qquad
\widehat\varphi\bigl(\unsure{\rho^{-1}}(\xi^{-1})\bigr)
=\unsure{\varepsilon_0},
\]
\[
\widehat\varphi\bigl(\rho^{-1}(\xi^{-1})\bigr)
\rho\bigl(\widehat\varphi(\rho^{-1}(\xi^{-1}))\bigr)^{-1}
=\unsure{\varepsilon_0\varepsilon_1^{-1}}
=(\sigma_\infty\rho)(\rho\sigma_\infty)^{-1}.
\]
}
\mspage{249}

\scanpage{251}
\mspage{426}

with
\[
(60)\qquad i=1,\quad\alpha_3=1\quad\text{i.e.}\quad
\varepsilon_3=-1,\quad j'=-1
\]

\uline{Remark} We can ask the question, in general,
of determining the elements
$u_{\alpha,\beta}=u_{\xi,\unsure{\eta}}$
of $\mathcal M_{0,3}^{!\sim}[0]$
(not necessarily normalized in one sense or the other)
such that we have (of course by hypothesis, $f=1$ ; but also)
\deleted{either} $\beta=1$ i.e. $\eta=1$,
\deleted{or} \deleted{\illegible{3}}
(\uline{N.B.} We already know that if $\alpha=1$ i.e.
$\xi=1$, then $u=1$)
We shall have to have
$\xi=\eta\ell_1^\nu=\ell_1^\nu$,
and it remains to check
\deleted{\illegible{2}}
\[
\xi\rho(\xi)
\rho^2(\xi)=1\quad\text{with}\quad
\xi=\ell_1^\nu,
\]
$\unsure{\mathrm{Sl}(2,\widehat{\mathbb Z})'}$,
\deleted{this} so that
\[
\underline{\xi}=
\begin{pmatrix}1&0\\2\nu&1\end{pmatrix}
\]
this means \inserted{(page \unsure{344|342})} that we have either
\[
\underline{\xi}=\pm\underline\rho^{-1}.
\quad\deleted{\text{i.e.}\ \underline{\xi}=\rho}
=\pm\begin{pmatrix}-1&1\\1&0\end{pmatrix}
\]
-- which is clearly not the case -- or
\[
C+D-B=\pm1\quad\text{i.e.}\quad2\nu+1=\mu=\pm1\quad\text{i.e.}
\]
$\nu=0$ or $-1$ -- these are the two trivial cases
$u_{\xi}=\mathrm{id}$ and
$u_{\xi}=\tau_\infty$.

\margin{We are \unsure{needlessly} repeating arguments already
made, we know very well that the multiplier is
(up to sign when the \unsure{normalization} of
$\xi$ has not been \unsure{done}) equal to one
of the ``\unsure{coefficients} \illegible{1}'' of
$\xi$, hence equal to $\pm1$.}
\mspage{250}

\scanpage{252}

Let us note \unsure{the|the}

\uline{Proposition.} The commutant of $\tau$ (the
nontrivial element of $\mathcal T_{0,3}^{!}\simeq\mathbb Z/2\mathbb Z
\subset\unsure{\mathcal N_{0,3}^{!\sim}|\mathcal M_{0,3}^{!\sim}}$)
in $\mathcal N_{0,3}^{!\sim}$ is reduced to $\{1,\tau\}$.

\margin{\uline{N.B.} This proposition \illegible{3}
in $\S\,46$ \unsure{VI|II}, \illegible{4} of \illegible{2}
of $\unsure{\widehat\Gamma}$.}

We shall do the calculation in the isomorphic group
$\unsure{\mathcal M_{0,3}^{!\sim}}(0)$,
\inserted{where $\tau$ corresponds to $\tau_\infty$},
in fact it turns out that the centralizer of $\tau$ in
the whole of $\mathcal M_{0,3}^{!\sim}[0]$
(the big \unsure{normalizer} of $L_0$ in
$\mathcal M_{0,3}^{!\sim}$, formed of the
$u_{\xi}$ without normalization conditions)
is reduced to $\{1,\tau\}$.
Indeed, we must express the commutation of a
$u=u_{\xi}$ with
$u_{\xi'}=\tau_\infty$,
where $\xi'=\ell_1^{-1}$, we have
\[
\left\{\begin{aligned}
u_{\xi'}u_{\xi}
&=u_{\xi''}
&&\text{with }\xi''=
\dot u_{\xi'}(\xi)\xi'
\\
u_{\xi}u_{\xi'}
&=u_{\xi'''}
&&\text{with }\xi'''=
\unsure{\dot u|u}_{\xi}(\xi')\xi
\end{aligned}\right.
\]
and we must express $\xi''=\xi'''$, i.e.
\[
(61)\qquad
\tau_\infty(\xi)\ell_1^{-1}
=u_{\xi}(\ell_1^{-1})\xi.
\]
Now we have
\[
\unsure{\dot u|u}_{\xi}(\varepsilon_1)
=\mathrm{int}(\xi)(\varepsilon_1^\mu)
\qquad\text{whence}
\]
\[
\unsure{\dot u|u}_{\xi}
\bigl(\underbrace{\ell_1^{-1}}_{\varepsilon_1^{-2}}\bigr)
=\mathrm{int}(\xi)(\ell_1^{-\mu})
\]
and (61) is written
\[
\tau_\infty(\xi)\ell_1^{-1}
=\xi\ell_1^{-\mu}\qquad\text{i.e.}
\]
\[
(62)\qquad
\xi^{-1}\tau_\infty(\xi)
=\ell_1^{1-\mu}=\ell_1^{-2\nu}.
\]
\mspage{251}

\scanpage{253}
\mspage{427}

This equation has the solution
\[
(63)\qquad\xi=\ell_1^\nu
\]
in $\widehat\Pi_{0,3}$ (for $\tau_\infty(\ell_1)=\ell_1^{-1}$
hence $\tau_\infty(\xi)=\ell_1^{-\nu}$)
and this solution is unique \uline{\unsure{if|once} we admit}
that the discrete results extend to the profinite case,
so as to have
\[
(64)\qquad\bigl(\widehat\Pi_{0,3}\bigr)^{\tau_\infty}=\{1\}.
\]
Now $\xi$ must moreover satisfy the relation
\[
\xi\rho(\xi)
\rho^2(\xi)=1
\]
and we have just seen that this implies
$\nu=\pm1$, i.e.
$\unsure{\dot u|u}\in\{\unsure{\mathrm{id}|1},\tau_\infty\}$.
\mspage{252}

\scanpage{254}
\mspage{428}

48. Construction of a group scheme.

I) Study of the map $\alpha\longmapsto
\alpha\rho(\alpha)^{-1}=[\alpha,\rho]$.

\unsure{Let|Let} $G$ be a group scheme over a ring
$A$ (for example $A=\mathbb Z$, $G=\mathrm{Sl}(2)_{\inserted{\mathbb Z}}
\deleted{\unsure{,\mathbb Z)}}$ or $\mathrm{Gl}(2)_{\mathbb Z}$),
$\rho$ a section of $G$, and let us consider the morphism
\[
(1)\qquad
F_\rho:\ \alpha\longmapsto\alpha\rho(\alpha)^{-1}
=[\alpha,\rho]:\qquad G\longrightarrow G.
\]
Let
\[
\deleted{\unsure{Z_\rho}}\quad
T_\rho
=\mathrm{Centr}_G(\rho)=G^\rho\subset G
\]
then
\[
(2)\qquad
F_\rho(\alpha')=F_\rho(\alpha)\quad\Longleftrightarrow\quad
\deleted{\unsure{\alpha}}\alpha'\deleted{\text{is}}
\text{ section of }\alpha T_\rho
\]
\inserted{thus} we have a canonical monomorphism
\inserted{induced by $F_\rho$}
\[
(3)\qquad
G/T_\rho\hookrightarrow G,
\]
whose image is to be characterized.

In the case that concerns us, the quotient group $G_{\mathrm{ab}}$
is representable (it is the trivial group for
$G=\mathrm{Sl}(2)_{\mathbb Z}$, and it is $\mathbb G_m$ if
$G=\mathrm{Gl}(2)_{\mathbb Z}$),
\deleted{and}\inserted{In any case,} a first condition for
a section $\xi$\inserted{$\in G(S)$}
over an $A$Scheme $S$, \deleted{over $A$} for it
to be in the image, is that we have
\[
(4)\qquad\delta(\xi)=1
\]
\mspage{253}

\scanpage{255}
\[
(5)\qquad \delta:G\longrightarrow G_{\mathrm{ab}}
\]
is the canonical \unsure{homomorphism}. A second condition is that $\xi\rho$, which must be (locally) of the form $\alpha\rho\alpha^{-1}$, be \uline{conjugate} \uline{locally} to $\rho$ -- and this condition is moreover necessary and sufficient, tautologically, for $\xi$ to belong to the image of $F_\rho$, and thus implies condition (4). But the latter may sometimes serve as an ingredient, \deleted{this} in a \illegible{1} of conditions of a ``computational'' nature, which are necessary and sufficient for us to have $\xi\rho\sim\rho$.

In the case of a reductive group $G$, the condition of \deleted{the \unsure{equivalence}}\inserted{\unsure{conjugation}} \inserted{locally} of sections is moreover fairly well \unsure{understood}, since in particular the work of Steinberg. We have a ``scheme of invariants'' $I_G$, and a canonical morphism
\[
(6)\qquad G\longrightarrow I_G
\]
\unsure{(}which is universal \inserted{(over $A$, and also after every base change)} for morphisms from $G$ into schemes, which are \unsure{invariant} under inner \unsure{automorphisms} -- i.e. $I_G$ is the \uline{representable} \uline{envelope} of $G/\mathrm{int}(G)$ \unsure{sheaf} of conjugacy
\mspage{254}

\scanpage{256}
\mspage{429}
classes.
\[
(7)\qquad G/\mathrm{int}(G)\longrightarrow I_G.
\]
To tell the truth, \deleted{it \illegible{1}} this morphism is an epimorphism (for \unsure{the|the} fppf topology), but not an isomorphism i.e. not a monomorphism -- in other words, for two sections (in this case $\rho$, $\xi\rho$) of $G$, it is indeed necessary, but not always sufficient, that their images in $I_G$ be equal, for them to be conjugate locally$^{t}$ fppf. But in the case where the two sections are \uline{\deleted{quasi}-regular} in the sense of Steinberg (which is a condition on each fiber) \deleted{it is}\inserted{the condition is} \unsure{sufficient}. It is also so if the two sections are \unsure{$\frac{1}{2}$} simple (which is also a condition on each fiber).

In the case of a group locally isomorphic to $\mathrm{Sl}(2)$, \inserted{to $\mathrm{GP}(1)$,} or $\mathrm{Gl}(2)$, the \uline{regular} sections in the sense of Steinberg are moreover those which are central at no point -- which is indeed the case for the $\rho$ or the $\sigma_\infty=\sigma$ of the preceding \S\. On the other hand, $I_G$ is then well known, \deleted{in the case of a group}
\mspage{255}

\scanpage{257}
\deleted{split, more generally of \uline{$\mathrm{Sl}(E)$} or \uline{$\mathrm{Gl}(E)$},}
$E$ locally free of rank 2, it is canonically isomorphic to $\mathbb E^1$ or to $\mathbb E^1\times\mathbb G_m$ (depending on whether the rank is 1 or 2) in the case of \deleted{\unsure{$\mathrm{Sl}(2)$|$\mathrm{Gl}(2)$}}a form of $\mathrm{Sl}(2)$ \deleted{\unsure{or $\mathrm{GP}(1)$}}, it is given by the \uline{\unsure{Trace}}
\[
(7)\qquad
\begin{array}{rcl}
G&\overset{\deleted{\mathrm{Tr}\;\unsure{\det}}}{\longrightarrow}&I_G\simeq\mathbb E^1,\\
x&\longmapsto&\mathrm{Tr}\,x
\end{array}
\]
in the case of a form of $\mathrm{Gl}(2)$ it is given by
\[
(8)\qquad
\begin{array}{rcl}
G&\longrightarrow&I_G=\mathbb E^{1*}\times\mathbb E^1.\\
x&\longmapsto&(\det x,\mathrm{tr}\,x)
\end{array}
\]
We thus find (for the record)

\uline{Proposition} Let $G$ be a form of $\mathrm{Sl}(2)$ \unsure{or} $\mathrm{Gl}(2)$ over a scheme $S$, let $\rho$ be a section of $G$ which is central in no fiber. For a section $\xi$ of $G$ to be \inserted{(locally)} in the image of $F_\rho$ ((1) or (3)), it is necessary and sufficient that it satisfy the following conditions

a) $\deleted{\unsure{\rho}}\xi\rho$ is central in no fiber, i.e. in no fiber do we have $\xi=\lambda\rho^{-1}$, where $\lambda$ \deleted{\unsure{is|being}} a central element

b) $\mathrm{tr}\,\xi\rho=\mathrm{tr}\,\rho$

c) (if $G$ a form of $\mathrm{Gl}(2)$) $\det\xi=1$ i.e. $\xi$ a section of $\mathrm{Sl}(2)$.

\margin{is equivalent to: $\det(\xi\rho)=\det\rho$}
\mspage{256}

\scanpage{258}
\mspage{430}
\uline{Corollary} When these conditions are satisfied, the scheme of solutions \inserted{in $\alpha$} of the equation
\[
[\alpha,\rho]=\xi
\]
is a (right) torsor under $T_\rho$, which is a smooth commutative group scheme over $S$, of relative dimension 1 (\unsure{if} $G$ a form of $\mathrm{Sl}(2)$) or 2 (if $G$ a form of $\unsure{\mathrm{Gl}(2)|\mathrm{Sl}(2)}$).

\uline{Examples \deleted{\unsure{1}}} $G=\mathrm{Gl}(2)_{\mathbb Z}$, $\rho=\begin{pmatrix}0&1\\-1&1\end{pmatrix}$ is the element of $\mathrm{Sl}(2,\mathbb Z)$ already so denoted (starting from page~335). It is indeed regular, \deleted{\unsure{and}} likewise the element $\sigma=\sigma_\infty=\begin{pmatrix}0&-1\\1&0\end{pmatrix}$.

In general, the datum of a form of $\mathrm{Gl}(2)$ is equivalent to that of a form of $\mathrm{Sl}(2)$, or of $\mathrm{GP}(1)$, or again to that of an Azumaya Algebra of rank 4 over $S$, let $M$ be one -- i.e. an Algebra locally \inserted{(fppf)} isomorphic to $M_2(\mathcal O_S)$ -- by
\[
(9)\qquad
\left\{
\begin{array}{rcl}
\unsure{\mathbb G|G}(M)&=&\text{scheme of units of }M\\
S\unsure{\mathbb G|G}(M)&=&\rule{4.8cm}{0.3pt}\quad\text{of determinant }1\\
P\unsure{\mathbb G|G}(M)&=&S\unsure{\mathbb G|G}(M)/\mathrm{Centre}\simeq\unsure{\mathbb G|G}(M)/\unsure{\mathbb G_m|G_m}
\end{array}
\right.
\]
\unsure{Thus|In} the datum \deleted{of $\rho$ defines} \unsure{corresponds} to that of an invertible section of $M$, and
\mspage{257}

\scanpage{259}
the regularity condition means that the subalgebra generated on each fiber is of rank 2 \inserted{(and not reduced to the scalars)} -- which implies that $\rho$ generates \deleted{\illegible{2}}\inserted{\unsure{a|unique} sub}-algebra of degree two
\[
(10)\qquad \mathcal O_S+\mathcal O_S\rho=A_\rho
\]
which is locally a direct \unsure{summand} in $M$ as a module, i.e. with formation \unsure{commuting|commuting} with base change \ldots\ And we shall have
\[
(11)\qquad
\begin{array}{l}
T_\rho=\text{scheme of units of }A_\rho\\
\qquad\text{(case of }G=\unsure{\mathbb G|G}(M)\text{)}
\end{array}
\]
or
\[
(11\text{bis})\qquad
\begin{array}{l}
T_\rho=\text{scheme of units of }A_\rho\\
\qquad\text{of \uline{norm} equal to }1\\
\qquad\text{(case of }G=S\unsure{\mathbb G|G}(M)\text{)}
\end{array}
\]
We shall note that $T_\rho$ is a \uline{torus} \inserted{(necessarily \unsure{maximal})} \unsure{iff} $\rho$ is $\frac{1}{2}$ simple, i.e. $\frac{1}{2}$ simple in each fiber -- which is equivalent, since it is regular, to the fact that in each fiber its two eigenvalues are distinct. For
\[
(12)\qquad \rho=\begin{pmatrix}0&1\\-1&\unsure{0|1}\end{pmatrix}\qquad (\rho^6=1,\ \rho^3=-1)
\]
this is \deleted{\unsure{is}} the case in characteristic $\ne 3$ -- \unsure{thus} \unsure{in the|in the} strict sense $T_\rho$ (\unsure{as} a group scheme over $\operatorname{Spec}\mathbb Z$) is \uline{not} a torus -- it
\mspage{258}

\scanpage{260}
\mspage{431}

becomes \deleted{\unsure{unipotent in characteristic 3}} \inserted{a product of a} unipotent with the center $\unsure{\mathbb G_m|G_m}$ in characteristic 3. Likewise, for
\[
\sigma=\sigma_\infty=\begin{pmatrix}0&-1\\1&0\end{pmatrix},\qquad
\sigma_{\unsure{\infty}}^4=1,\quad\sigma^2=-1,
\tag{13}
\]
$\sigma$ is $\unsure{\frac12}$ simple in all characteristics except in characteristic 2,
where it is unipotent -- and $T_\sigma$ is a
torus except in characteristic 2, where it \deleted{becomes} is
a product of the \inserted{center} $\unsure{\mathbb G_m|G_m}$ with a unipotent.

\margin{Let us return to the general case}
When we take the centralizer in
$S\unsure{\mathbb G|G}(M)$, we find
\[
ST_\rho\subset T_\rho,\qquad
ST_\rho=\operatorname{Ker}\left(
T_\rho
\underset{\inserted{\substack{\text{\unsure{is induced|induced}}\\\text{by the determinant}}}}
{\xrightarrow{\text{norm}}}\unsure{\mathbb G_m|G_m}\right)
\tag{13}
\]

The advantage of the $T_\rho$ over $ST_\rho$ is that
every torsor under $T_\rho=\text{scheme of units of }A_\rho$
is locally trivial in the Zariski sense, whereas
this is no longer necessarily so for
a torsor under $ST_\rho$. Thus, when
we set out to solve the equation in $\alpha$
\[
[\alpha,\rho]=\xi,
\]
for given $\xi$ satisfying the conditions of the
proposition, this is possible (Zariski locally) in
$\unsure{\mathbb G|G}(M)$, but perhaps not in $S\unsure{\mathbb G|G}(M)$.

\mspage{259}

\scanpage{261}

Let us return to the special case
\[
\rho=\begin{pmatrix}0&1\\-1&1\end{pmatrix},\qquad A=\mathbb Z
\]
then
\[
\begin{gathered}
A_\rho=\text{ring of integers of the field}\\
\mathbb Q(\sqrt[3]{-1}),\quad \rho\text{ corresponds to}\\
\exp(2i\pi/6)=\frac{1+i\sqrt3}{2}
\end{gathered}
\tag{14}
\]

When we give \unsure{ourselves} $\xi$, \deleted{the conditions of the}
\[
\xi=\begin{pmatrix}A&B\\C&D\end{pmatrix},\quad
\text{whence }\xi\rho=
\begin{pmatrix}-B&A+B\\-D&C+D\end{pmatrix}
\LMcancelledtag{\deleted{\unsure{15}}}
\]
the conditions \deleted{\unsure{of}} b), c) of the proposition
\deleted{become}
\[
\det\xi=1,\quad \operatorname{Tr}(\xi\rho)=\operatorname{Tr}\rho
\]
\deleted{(16)} become
\[
AD-BC=1\qquad C+D-B=1
\tag{16}
\]
-- quite familiar conditions! They imply
the possibility \deleted{of \unsure{finding}} (over an arbitrary scheme)
of finding Zariski locally an $\alpha\in\Gamma\mathrm{Gl}(2)$
such that \deleted{\unsure{satisfies}}
\[
\alpha\rho(\alpha)^{-1}=[\alpha,\rho]=\xi,
\tag{\unsure{*}}
\]
\margin{at least if we suppose that the condition \unsure{of regularity} is not \unsure{violated} on any fiber\ldots}
and over $\mathbb Z$ it is even possible globally,
for the obstruction is an \deleted{\unsure{torsor under}}
\inserted{invertible $A_\rho$-Module, now}
$A_\rho^{\deleted{*}}$ is principal \unsure{unless I am mistaken} i.e. \deleted{\unsure{Pic of}}
\[
\operatorname{Pic}(A_\rho)=1
\tag{17}
\]

\mspage{260}

\scanpage{262}
\mspage{432}

But if $\alpha$ is a solution, we shall have a priori
\[
\det\alpha\in\mathbb Z^*=\{\pm1\}
\tag{18}
\]
and \deleted{every} other solution will be of the form
$\alpha'=\alpha f$, with $f\in A_\rho^*$, whence
\[
\det\alpha'=\det\alpha\cdot\det f=\det\alpha
\]
for
\[
\det f=f\overline f>0\quad\text{hence }\det f=+1,
\]
thus the sign in (18) \deleted{\unsure{is}} does not depend
on the choice of $\alpha$ -- if it is $-1$, one cannot
find a solution of equation (*) in
$\mathrm{Sl}(2,\mathbb Z)$. The sign depends only on $\xi$ i.e. only
on $A,B,C,D$ -- and in fact \deleted{\unsure{if $D$}} only on $B,D$
(which already determine $\xi$, thanks to
formulas (16) \inserted{at least if $D$ \unsure{invertible}}) -- the computations made elsewhere
suggest that this sign is also \deleted{\illegible{1}} equal to the
remainder of $D$ mod 3 -- $1$ if $D\equiv1\ (3)$, $-1$ if
$D\equiv-1\ (3)$.

Returning to the general case, if we have found
a relation
\[
[\alpha,\rho]=\alpha\rho(\alpha)^{-1}=\xi
\qquad\alpha\in\Gamma\unsure{\mathbb G|G}(M),
\tag{19}
\]
in order \deleted{\unsure{that there exist}} to find a solution $\alpha'$ which
is $\in\Gamma S\unsure{\mathbb G|G}(M)$ i.e. $\det\alpha'=1$, we know that
the solutions $\alpha'\in\Gamma\unsure{\mathbb G|G}(M)$ \deleted{\unsure{sts}} are given by
\[
\alpha'=\alpha f\qquad f\in\Gamma\underline{\mathcal O}_S^*,
\]

\mspage{261}

\scanpage{263}

and the condition becomes
\[
\det\alpha'=1\quad\text{i.e.}\quad\det(\alpha f)=1\quad\text{i.e.}
\]
\[
\det(f^{-1})=\det\alpha
\]
thus the necessary and sufficient condition for the existence
of a solution is that the section $\det\alpha$ of
$\mathcal O_S^*$ be in the image of the norm
homomorphism
\[
\unsure{N|\mathcal N}_{A_\rho/A}:A_\rho^*\longrightarrow A^*,
\tag{20}
\]
\margin{N.B. $S=\operatorname{Spec}A$}
i.e. the obstruction is in the group
\[
A^*/\unsure{N|\mathcal N}_{A_\rho/A}A_\rho^*
\tag{21}
\]
which is a particularly well-known invariant
in class field theory!

\margin{\uline{N.B.} $N(c\rho+d)=\unsure{c^2+cd+d^2}$}
Unless I am mistaken, it follows from class field
theory that for $\rho$ given by (12),
and for $A=\underbrace{\mathbb Z_\ell}_{\inserted{\text{integers }\ell\text{-adic}}}$ ($\ell$ a prime number), this
quotient group is zero in all cases,
except for $\deleted{\unsure{\rho|\ell}}\ \ell=3$, where it is isomorphic to
$\mathbb Z/2\mathbb Z$. \deleted{\illegible{1} Let us note}

\margin{arbitrary \unsure{base|Ring}}
Let us note that (still in the special case (12))
if
\[
\alpha=\begin{pmatrix}a&b\\c&d\end{pmatrix}
\tag{22}
\]
immediate computations already made show that

\mspage{262}
